\setcounter{page}{145}
\section{\S2 The Functor $f^!$ for a Smooth Morphism}

\textbf{Definition.}
Let $f\colon X\to Y$ be a smooth morphism of preschemes.
Define the functor
\[
f^!\colon D(Y)\to D(X)
\]
by
\[
f^\bullet(G^\bullet) = \mathbf{L}f^*\big(G^\bullet\big) \otimes \omega_{X/Y}[n],
\]
where $[n]$ denotes the shift left by $n$ degrees.

Since smooth morphisms are flat, we have $\mathbf{L}f^* = f^*$ on all of $D(Y)$.
The dualizing sheaf $\omega_{X/Y}$ is invertible, hence belongs to $D_{\mathrm{fTd}}(X)$,
so the derived tensor product is well-defined (cf.\ Chapter II).

\textbf{Proposition 2.1.}
Let $f\colon X\to Y$ be smooth, $u\colon Y'\to Y$ a morphism of finite Tor-dimension,
and let $X' = X\times_Y$. Then there is a natural isomorphism of functors
\[
\mathbf{L}v^*\circ f^! \;\cong\; g^!\circ \mathbf{L}u^*
\]
from $D(Y)$ to $D(X')$.

\textit{Proof.}
Follows from the base-change compatibility of $\omega_{X/Y}$ (\S1)
and derived functor compatibilities in Chapter II.

\setcounter{page}{146}
\textbf{Remark.}
We follow the convention of writing "$=$" for canonical isomorphisms whenever the compatibility is unambiguous.
When sign conventions or coordinate choices affect the isomorphism, we explicitly name the morphism and track compatibilities.

\textbf{Proposition 2.2.}
Let $X\xrightarrow{f}Y\xrightarrow{g}Z$ be smooth morphisms.
There exists a functorial isomorphism
\[
\zeta_{f,g}\colon (gf)^! \;\xrightarrow{\sim}\; f^!\circ g^!
\]
of functors $D(Z)\to D(X)$.
For any composition of three smooth morphisms, the corresponding isomorphisms yield commutative diagrams.

\textit{Proof.}
Construct $\zeta_{f,g}$ using the composition isomorphism for dualizing sheaves (Definition 1.5).
Commutativity follows from Proposition 1.6.

\textbf{Proposition 2.3.}
Let $f\colon X\to Y$ be smooth.
Then $f^!$ maps $D_{\mathrm{qc}}(Y)$ to $D_{\mathrm{qc}}(X)$.
If $X,Y$ are locally noetherian, $f^!$ maps $D_{\mathrm{c}}(Y)$ to $D_{\mathrm{c}}(X)$.

\textit{Proof.}
Immediate from quasi-coherence and coherence pullback properties.

\setcounter{page}{147}
\textbf{Proposition 2.4.}
Let $f\colon X\to Y$ be smooth.
\begin{enumerate}
\item[(a)] There is a functorial isomorphism
\[
f^\bullet\big(F^\bullet \stackrel{\mathbf{L}}{\otimes} G^\bullet\big)
\;\cong\;
f^\bullet F^\bullet \stackrel{\mathbf{L}}{\otimes} f^\bullet G^\bullet,
\]
valid if either $F^\bullet,G^\bullet\in D^-(Y)$,
or one lies in $D^b_{\mathrm{fTd}}(Y)$ and the other in $D(Y)$.

\item[(b)] There is a functorial homomorphism
\[
f^\bullet\big(\mathbf{R}\mathcal{H}\textit{om}^\bullet(F^\bullet,G^\bullet)\big)
\to
\mathbf{R}\mathcal{H}\textit{om}^\bullet\big(f^\bullet F^\bullet,\,f^\bullet G^\bullet\big),
\]
for $F^\bullet\in D(Y)$, $G^\bullet\in D^+(Y)$.
If $Y$ is locally noetherian and $F^\bullet\in D_{\mathrm{c}}^-(Y)$, this is an isomorphism.
\end{enumerate}

\textit{Proof.}
Follows from Chapter II Propositions 5.8, 5.9, 5.13, 5.16.